\begin{abstract}
Diffusion models generate samples by traversing a denoising trajectory, a sequence of noise-reduction steps that progressively transforms random noise into data of the target distribution. A practical question at inference time is how to use a given compute budget to generate higher-quality samples without retraining. Standard approaches allocate compute uniformly across the denoising steps, but this is suboptimal: early steps that determine global structure often benefit more from additional refinement than late steps that refine local details. We formalize this as a budget allocation problem: at each timestep $t$, evaluate $K_t$ noise candidates and keep the best (under a fixed quality scorer, the \emph{verifier}); choose $\{K_t\}_{t=1}^T$ to maximize the expected verifier score of the final sample subject to $\sum_t K_t = B$ and $K_t \ge 1$. We develop a theoretical model for this problem and establish two optimality results. First, the optimal offline allocation has a closed-form water-filling rule: each step's share grows with its sensitivity to noise refinement. Second, we prove an $\Omega(T)$ worst-case regret lower bound against any adaptive policy, motivating an offline-plus-online design rather than pure online adaptation. We instantiate this analysis as \textsc{GAINS} (Global Adaptive Inference-time Noise Scheduling), an algorithm that pairs offline sensitivity profiling with online feedback control. On Stable Diffusion, EDM, and flow-based models, \textsc{GAINS} matches the quality of uniform allocation using 20--50\% fewer function evaluations, and exceeds it at equal budget.
\end{abstract}
